
\clearpage
\VerdierPhysicalPageAnchor{0129}%
\begingroup
\small
\setlength{\emergencystretch}{1.5em}%
\setlength{\abovedisplayskip}{3pt plus 2pt minus 2pt}%
\setlength{\belowdisplayskip}{3pt plus 2pt minus 2pt}%

\VerdierSemanticAnchor{ch2-2-1-1-sm3-conditions}
\VerdierSemanticAnchor{ch2-2-1-1-sm3-i}
\noindent\textit{i)} 始域が $Y$ である $S$ の射 $s$ が存在し、$sf=sg$ となる。

\VerdierSemanticAnchor{ch2-2-1-1-sm3-ii}
\noindent\textit{ii)} 終域が $X$ である $S$ の射 $t$ が存在し、$ft=gt$ となる。

\VerdierSemanticAnchor{ch2-2-1-1-saturated}
\noindent 乗法的系が次の性質をもつとき、\emph{飽和}であるという。

\VerdierSemanticAnchor{ch2-2-1-1-sm4}
\VerdierSemanticAnchor{ch2-2-1-1-saturated-sm4}
\smallskip
\noindent\textbf{SM4)} $f,g,$ と $h$ が合成可能な三つの射で、$fg$ と $gh$
が $S$ に属するならば、$g$ は $S$ に属する。

\VerdierSemanticAnchor{ch2-2-1-2}
\medskip
\noindent\textbf{2.1.2.}\quad
\VerdierSemanticAnchor{ch2-2-1-2-compatible-triangulation}
\VerdierSemanticAnchor{ch2-2-1-2-compatible}
$\mathcal D$ を三角圏とする。$\mathcal D$ の乗法的系 $S$ が次の性質を
もつとき、\emph{三角構造と両立する}という。

\VerdierSemanticAnchor{ch2-2-1-2-sm5}
\smallskip
\noindent\textbf{SM5)} $S$ の各元 $s$ と各整数 $n\in\mathbb Z$ に対し、
射 $s[n]$ は $S$ に属する。言い換えれば、$S$ は\emph{シフト自己同型の作用で
安定}である。

\VerdierSemanticAnchor{ch2-2-1-2-sm6}
\smallskip
\noindent\textbf{SM6)} 任意の二つの完全三角

\VerdierSemanticAnchor{ch2-2-1-2-sm6-triangle-pair}
\[
  (X,Y,Z,u,v,w), \qquad (X',Y',Z',u',v',w')
\]
\noindent および任意の可換図式

\VerdierSemanticAnchor{ch2-2-1-2-sm6-square}
\begin{center}
  \begin{tikzcd}[column sep=4.2em,row sep=2.5em]
    X \arrow[r,"u"] \arrow[d,"s"'] & Y \arrow[d,"s'"] \\
    X' \arrow[r,"u'"'] & Y'
  \end{tikzcd}
  \quad ,
\end{center}
\VerdierSemanticAnchor{ch2-2-1-2-sm6-triangle-morphism}
\VerdierSemanticAnchor{ch2-2-1-2-sm6-completion}
\noindent において、$s$ と $s'$ が $S$ に属するならば、$S$ に属する射
$s''\colon Z\to Z'$ が存在し、$(s,s',s'')$ は三角の射となる。

\VerdierSemanticAnchor{ch2-2-1-3}
\medskip
\noindent\textbf{注 2.1.3.}\quad
\VerdierSemanticAnchor{ch2-2-1-3-item1}
\VerdierSemanticAnchor{ch2-2-1-3-1}
\textbf{1)} $\mathcal D$ を三角圏、$S$ を三角構造と両立する
$\mathcal D$ の乗法的系、$\mathcal D^\circ$ を $\mathcal D$ の反対三角圏
とする。$\mathcal D^\circ$ の射の集合としての $S$ は、$\mathcal D^\circ$
の乗法的系であり、$\mathcal D^\circ$ の三角構造と両立する (1.1.7)。

\VerdierSemanticAnchor{ch2-2-1-3-item2}
\VerdierSemanticAnchor{ch2-2-1-3-2}
\smallskip
\noindent\textbf{2)} $\mathcal D$ を三角圏とし、$S$ を性質 SM1)、SM5)、
SM6) をもつ $\mathcal D$ の射の集合とする。このとき $S$ は性質 SM2) も
もつ。
\VerdierSemanticAnchor{ch2-2-1-3-2-sm2}
実際、任意の図式

\VerdierSemanticAnchor{ch2-2-1-3-item2-input-boundary}
\VerdierSemanticAnchor{ch2-2-1-3-2-input-boundary}
\begin{center}
  \begin{tikzcd}[column sep=4.0em,row sep=2.6em]
    & X' \arrow[d,"s\in S"] \\
    Z \arrow[r,"f"] & X
  \end{tikzcd}
\end{center}
\endgroup
